\subsection{罗尔定理}

命题 1（“罗尔定理”）. 设 f 是实函数，在闭区间 I = {[}a, b{]} （其中 a \textless{} b）上有限且连续，在 {]}a, b{[} 的每一点具有导数（有限或无穷），且满足 \(f(a) = f(b)\) 。则存在点 c ∈ {]}a, b{[} ，使 \(f'(c) = 0\) 。

若 \(f\) 为常数，则命题是显然的；若不然，则 \(f\) 例如取 \(> f(a)\) 的值，于是在 I 的一个内点 \(c\) 处达到它的上确界（Gen.~Top., IV, p.~359, th. 1）。由于 \(f\) 在该点具有相对极大值，我们有 \(f'(c) = 0\) （I, p.~20, prop. 7）。

推论. 设 \(f\) 是实函数，在 \([a, b]\) （其中 \(a < b\) ）上有限且连续，并在每一点具有导数（有限或无穷）。则存在点 \(c\) ∈ \(]a\) , \(b[\) ，使 \(f(b) - f(a) = f'(c)(b - a)\) 。

只需对函数 \(f(x) - \frac{f(b) - f(a)}{b - a} (x - a)\) 应用命题 1。

这一推论表明，在 f 的图像 C 上存在一点 \(\mathbf{M}_{c} = (c, f(c))\) ，使得 a \textless{} c \textless{} b ，且 C 在该点的切线平行于连接点 \(\mathbf{M}_{a} = (a, f(a))\) 与 \(\mathbf{M}_{b} = (b, f(b))\) 的直线。

